\begin{frame}{17.3 확인 문제}
  \small
  \begin{enumerate}
    \item SFT 손실 $\mathcal{L} = -\frac{1}{T}\sum
      \log \pi_\theta(y_t | y_{<t})$에서, SFT의 \textbf{입력
      데이터}는 (질문, \textbf{정답 답변}) 쌍 — 이 ``정답
      답변''을 누가 작성하는가?
    \item RLHF 2단계에서 \textbf{보상모델}이 학습받는
      손식은? \textbf{정책}이 학습받는 목적함수(17.3의
      $J(\theta)$)를 쓰고, KL 항이 \textbf{왜} 필요한가?
    \item DPO 손실을 BT 손실과 \textbf{비교}: BT의
      $\log r(x, y_w) - \log r(x, y_l)$ 대신 들어가는
      항을 쓰고, $\beta$가 어떤 역할을 하는가.
    \item \textbf{판별:} ``DPO는 보상모델이 필요 없으므로
      RLHF보다 항상 안정적이다'' — 참/거짓, 이유(보상 해킹
      관점).
    \item LoRA rank $r = 8$을 $r = 64$로 올리면:
      (a) 학습 파라미터 수가 몇 배? (b) \textbf{과적합}
      위험은 어떻게 변하는가? (c) 17.3의 SVD 근사
      결과(rank-4: 0.0\%, rank-3: 0.3\%, rank-1: 0.7\%)
      와 비교.
    \item RAG: 질문과 \textbf{검색된 문서}를 컨텍스트에
      넣는다 — 17.2의 ``컨텍스트 = \textbf{조건}'' 관점에서
      RAG는 \textbf{학습}인가, \textbf{조건화}인가?
  \end{enumerate}
\end{frame}
